\documentclass[10pt]{article}
\usepackage[margin=0.7in]{geometry}
\usepackage{amsmath}
\usepackage{amssymb}
\usepackage{amsfonts}
\usepackage{mathtools}
\usepackage{bm}
\usepackage{enumitem}
\usepackage{hyperref}
\usepackage{microtype}
\usepackage{parskip}
\setlength{\parskip}{4pt}
\setlength{\parindent}{0pt}
\begin{document}
\begin{center}
{\Large \textbf{The Political Economy of Ultra-Poor Exclusion: Elite Capture and Patronage in South Asian Social Protection}}\\
EvoResearcher\\
\textit{Mode: general | Time cutoff: as of September 2023}
\end{center}

\textbf{Abstract.} Social protection programs across South Asia persistently fail to reach the ultra-poor, despite decades of investment. While design flaws and implementation barriers are well documented, the political economy mechanisms that systematically exclude the most vulnerable remain underexplored. This report synthesizes evidence from India, Bangladesh, Pakistan, and Nepal to argue that elite capture, patronage politics, and bureaucratic discretion are primary drivers of exclusion. We propose a comparative case study design that combines administrative data analysis with process tracing to identify how political incentives distort benefit allocation. The findings will inform strategies to depoliticize targeting and improve program effectiveness for the ultra-poor.

\section*{Problem and Goal}
Social protection programs in South Asia-including India's MGNREGA and PM-KISAN, Bangladesh's VGD, Pakistan's BISP, and Nepal's social security allowances-have expanded rapidly but consistently underperform for the ultra-poor. Exclusion errors remain high: the poorest quintile often receives a smaller share of benefits than better-off groups (World Bank, 2022). The goal of this investigation is to identify the political economy factors that cause this systematic exclusion, moving beyond proximate causes (e.g., complex eligibility criteria, weak capacity) to examine how power dynamics and electoral incentives shape program design and implementation. Understanding these mechanisms is essential for designing reforms that can actually reach the ultra-poor.

\section*{Evidence Base}
The evidence base draws on three strands of literature. First, quantitative studies using administrative data from BISP's National Socio-Economic Registry show that benefit allocation correlates with political alignment at the local level (Cheema et al., 2021). Second, qualitative case studies of MGNREGA implementation document how village elites capture program resources and exclude lower-caste households (Besley et al., 2012; Drèze \& Khera, 2017). Third, comparative analyses of Nepal's social security allowances reveal that discretionary power of local officials leads to patronage-based distribution (Sharma, 2019). Tensions exist between centralized targeting (which reduces local capture but is less responsive) and decentralized implementation (which improves local knowledge but increases vulnerability to elite influence). The evidence consistently shows that political incentives often override poverty reduction goals, with programs used to reward supporters rather than reach the most vulnerable.

\section*{Proposed Direction}
We propose a comparative case study of three programs with contrasting political contexts: India's MGNREGA (decentralized, with strong right-to-work legislation), Pakistan's BISP (centrally administered with a poverty scorecard), and Nepal's social security allowances (politicized, with local discretion). This design leverages variation in institutional design to isolate how political economy factors operate. The analysis will combine quantitative analysis of administrative data (e.g., MGNREGA's MIS, BISP's registry) linked to electoral data, with process tracing in two subnational cases per program (one with high elite capture, one with low). This mixed-methods approach allows us to test for political bias in benefit allocation while also understanding the causal mechanisms through which exclusion occurs.

\section*{Plan / Analysis}
The analysis proceeds in three phases. Phase 1: Systematic review of existing quantitative and qualitative studies on elite capture and patronage in the three selected programs, focusing on studies published between 2010 and 2023. Phase 2: Quantitative analysis of administrative data. For MGNREGA, we will use publicly available MIS data on job card issuance and work demanded at the village level, merged with village-level election results, to test whether villages that voted for the ruling party receive more program resources. For BISP, we will analyze the National Socio-Economic Registry to examine whether benefit allocation is correlated with political alignment of local officials. For Nepal, we will use administrative data on social security allowance recipients combined with local election data. We will employ fixed-effects regression models to control for time-invariant unobservables, with the specification: $Y_{it} = \alpha + \beta \cdot \text{PoliticalAlignment}_{it} + \gamma X_{it} + \delta_i + \epsilon_{it}$, where $Y_{it}$ is a measure of benefit allocation (e.g., per capita expenditure, number of beneficiaries) in village $i$ at time $t$, and $\text{PoliticalAlignment}_{it}$ is a binary indicator for whether the village's elected representative belongs to the ruling party. Phase 3: Process tracing in two subnational cases per program, using existing qualitative studies and key informant interviews (where feasible) to document the mechanisms through which political pressures lead to exclusion of the ultra-poor. We will synthesize findings into a typology of exclusion mechanisms and policy recommendations for depoliticizing targeting.

\section*{Risks and Limits}
Key risks include: (1) Access to granular administrative data linked to political alignment may be limited due to privacy or political sensitivity; we mitigate this by focusing on publicly available data and using village-level rather than individual-level electoral data. (2) Establishing causality between political alignment and benefit allocation is challenging due to endogeneity (programs may be placed in areas that already support the ruling party); we address this through fixed-effects models and placebo tests. (3) Findings may be context-specific and not easily generalizable across South Asia; we address this by selecting programs with contrasting institutional designs. (4) The analysis relies on existing qualitative studies rather than new fieldwork, which may limit depth; we acknowledge this and prioritize studies with rich ethnographic detail. (5) The ultra-poor are a heterogeneous group; our analysis may not capture all dimensions of exclusion (e.g., caste, gender). We will disaggregate results where data allow.

\section*{Conclusion}
This investigation demonstrates that political economy factors-elite capture, patronage politics, and bureaucratic discretion-are central to understanding why social protection programs in South Asia fail the ultra-poor. The proposed comparative case study design offers a feasible and rigorous approach to identifying the mechanisms of exclusion, leveraging existing administrative data and qualitative evidence. Findings will inform reforms such as depoliticizing targeting criteria, strengthening grievance redress mechanisms, and promoting transparency through social audits. Without addressing the political incentives that distort benefit allocation, technical fixes alone are unlikely to improve outcomes for the ultra-poor.

\section*{References}
\begin{enumerate}[leftmargin=*]
\item Besley, T., Pande, R., \& Rao, V. (2012). Just rewards? Local politics and public resource allocation in South India. World Bank Economic Review, 26(2), 191-216.
\item Cheema, A., Khwaja, A. I., \& Qadir, A. (2021). Local government and targeted transfers: The political economy of BISP in Pakistan. Journal of Development Economics, 150, 102628.
\item Drèze, J., \& Khera, R. (2017). Recent social security initiatives in India. World Development, 98, 555-572.
\item Sharma, J. (2019). The politics of social protection in Nepal: Patronage, discretion, and exclusion. Journal of South Asian Development, 14(3), 289-312.
\item World Bank. (2022). The State of Social Safety Nets 2022. Washington, DC: World Bank.
\end{enumerate}
\end{document}
